\section{Implementation: layout, CI, replay}
\label{sec:appendix-implementation}

The pattern is materialised in the public repository that produced
this paper.\footnote{\url{https://github.com/noheton/f-ai-r}} The
intent is not to prescribe a stack: any version-control system with
a build target, a structured-data format that supports PROV-O, and a
way to version prompt files would do.

\paragraph{Layout.}
\texttt{paper/} carries the \LaTeX{} sources, bibliography, Makefile,
and figure sources, with one \texttt{\textbackslash section\{\}}
per file under \texttt{paper/sections/}. \texttt{doc/} carries the
methodology, FAIR notes, human-AI collaboration policy, research
protocol, submission plan, the append-only logbook, the PROV-O graph
in Turtle, and a Mermaid mirror. \texttt{agents/} carries one
markdown file per role (orchestrator, scientific writer, source
analyzer, FAIR aligner, provenance curator, layout scrutinizer,
readability reviewer, illustration, condenser, research protocol),
each structured per the prompt taxonomy of
Appendix~\ref{app:prompts}. \texttt{scripts/} carries the static-site
generator that parses \texttt{doc/provenance.ttl} via \texttt{rdflib}
and renders it into a small HTML site.

\paragraph{Continuous integration.}
\texttt{build-paper.yml} compiles the manuscript on every push that
touches \texttt{paper/} and uploads the PDF as an artefact, using
\texttt{xu-cheng/latex-action@v3}. \texttt{pages.yml} validates the
Turtle, builds the site, and (on \texttt{main}) deploys via the
official GitHub Pages action chain. Working branches see the build
output as a downloadable artefact, never as a public deploy.

\paragraph{Replay loop.}
A reader who clones the repository, installs a TeX Live distribution
and the dependencies in \texttt{scripts/requirements.txt}, runs
\texttt{make -C paper pdf} and
\texttt{python scripts/build\_provenance\_site.py}, has reproduced
the paper, the provenance graph, and the public site. The audit pass
for \textbf{Reusable} is, in our setup, that loop run on a clean
machine, compared against the published artefact. Any difference is
a defect.
